% Generated by tools/edn_latex.py from reports/edn.md.
% Do not edit: change reports/edn.md and run `make edn`.
\ednentry{
  id = {EDN-09},
  title = {Bump to Python 3.12, to use real SHAP instead of a workaround},
  date = {2026-09-22},
  milestone = {M2: Reproducibility},
  topic = {Tooling, Explainability},
  participants = {@lukas2510},
  decision = {\texttt{requires-\allowbreak{}python} is \texttt{\textasciitilde{}=\allowbreak{}3.\allowbreak{}12.\allowbreak{}0} (was \texttt{\textasciitilde{}=\allowbreak{}3.\allowbreak{}11.\allowbreak{}0}). \texttt{pyproject.\allowbreak{}toml}, \texttt{uv.\allowbreak{}lock}, and the Python version in \texttt{.\allowbreak{}github/\allowbreak{}workflows/\allowbreak{}ci.\allowbreak{}yml} are updated, and \texttt{shap} becomes a usable dependency for offline analysis. Which SHAP implementation FR-08's explanation endpoint uses is decided separately in EDN-11, which amends this entry: the API computes contributions from the booster's own export, and \texttt{shap} stays out of the API image.},
  rationale = {\emph{Alternatives considered:}
\begin{itemize}[leftmargin=*, topsep=2pt]
\item \textbf{Option A: use LightGBM's native \texttt{predict(pred\_\allowbreak{}contrib=\allowbreak{}True)} instead of the \texttt{shap} package.} Gives the same exact TreeSHAP values without touching \texttt{shap} or its Python 3.12 requirement, and avoids \texttt{numba}/\texttt{llvmlite} in the API image.
\newline Pros: no project-wide version bump needed this early in the semester; smaller API image.
\newline Cons: reimplements what \texttt{shap} already provides; still needs the real \texttt{shap} package for any offline/notebook analysis (global importance, summary plots), so the Python constraint would resurface there anyway; more code to maintain for a workaround to a constraint that turned out to be removable.
\item \textbf{Option B (chosen): bump to Python 3.12 project-wide.} Verified empirically: \texttt{uv lock -\allowbreak{}-\allowbreak{}python 3.\allowbreak{}12} and \texttt{uv sync -\allowbreak{}-\allowbreak{}python 3.\allowbreak{}12} resolved and installed cleanly with the full planned M2-M5 dependency set added (\texttt{shap}, \texttt{mlflow>=\allowbreak{}2.\allowbreak{}22,<3}, \texttt{lightgbm}, \texttt{catboost}, \texttt{mapie}, \texttt{fastapi}, \texttt{uvicorn}, \texttt{pydantic}, \texttt{great-\allowbreak{}expectations}, \texttt{dvc}, \texttt{pytest-\allowbreak{}cov}, \texttt{codecarbon}, \texttt{httpx}), and \texttt{shap.\allowbreak{}TreeExplainer} on a trained LightGBM model produced correct SHAP values under 3.12 in a functional test.
\newline Pros: uses \texttt{shap} directly with no workaround; every other planned dependency (checked against PyPI's \texttt{requires\_\allowbreak{}python} metadata for the exact pinned versions, not just ``latest'') only has a lower bound compatible with 3.12, so nothing else in the stack is put at risk; no Dockerfile exists yet to update, only three lines in \texttt{ci.\allowbreak{}yml}.
\newline Cons: a project-wide version bump this early, touching \texttt{pyproject.\allowbreak{}toml}, \texttt{uv.\allowbreak{}lock} and CI, for the sake of one requirement.
\end{itemize}
\par
\emph{Why:} The constraint was checked empirically rather than assumed from documentation: a scratch resolve of the full planned dependency set under Python 3.12 succeeded with no conflicts, and \texttt{shap.\allowbreak{}TreeExplainer} functionally worked. With no blocker found, option B removes the constraint at its source instead of working around it, and the change is small (no Dockerfile yet, three \texttt{ci.\allowbreak{}yml} lines, \texttt{pyproject.\allowbreak{}toml}, \texttt{uv.\allowbreak{}lock}).},
  ai = {Information seeking, Alternative generation, Alternative assessment, Recommendation},
  response = {Accepted},
  assessment = {While reviewing PR \#19, AI proposed option A as a way around the Python 3.11/SHAP 0.52 conflict documented in the brief. Asked to check whether the version could be bumped instead, AI verified this empirically (a scratch \texttt{uv lock}/\texttt{uv sync} under 3.12 with the full planned dependency set, plus a functional \texttt{TreeExplainer} test) rather than relying on PyPI metadata alone, found no blocker, and presented both options with the empirical evidence. Lukas reviewed it and chose to bump.},
  aievidence = {Claude Code session on 2026-09-22 while reviewing PR \#19: AI first recommended \texttt{pred\_\allowbreak{}contrib} as a workaround, Lukas asked ``please check if we can bump our python version,'' AI ran an empirical \texttt{uv lock}/\texttt{uv sync}/import/functional check under Python 3.12 with the full planned stack, found no blocker, and reported the evidence; Lukas replied ``ok please then bump the version in the PR an adjust the requriemnt.''},
  evidence = {\href{https://github.com/mlops-2627q1-mds-upc/MLOPS_Recommenditos/blob/main/docs/docs/specification.md}{Specification} FR-08; \href{https://github.com/mlops-2627q1-mds-upc/MLOPS_Recommenditos/blob/main/docs/docs/project-brief.md}{project brief} section 6 (tooling constraints); \texttt{pyproject.\allowbreak{}toml}, \texttt{.\allowbreak{}github/\allowbreak{}workflows/\allowbreak{}ci.\allowbreak{}yml}; \href{https://github.com/mlops-2627q1-mds-upc/MLOPS_Recommenditos/pull/19}{PR \#19}.},
}
